\chapter*{Abstract}
\addcontentsline{toc}{chapter}{Abstract}
\begin{english}

The convergence of terrestrial 5G/6G networks (TN) with Low Earth Orbit
(LEO) satellite networks into Satellite-Terrestrial Integrated Networks
(STINs) is a core design goal of 6G, aimed at extending coverage to areas
where terrestrial infrastructure is limited or absent. Multi-Connectivity
(MC) between base stations and satellites has been proposed as a way to
improve reliability, coverage, and load management in such networks.

This thesis develops a MATLAB simulation framework for a connectivity
scenario consisting of two terrestrial base stations (5G NR, 3GPP TR
38.901 UMa/UMi) and one LEO satellite, in which each user connects to
every candidate node (terrestrial, satellite) whose SNR clears a minimum
usable threshold -- simultaneously to both when both clear it
(multi-connectivity, greedy per-user decision). The model computes
propagation loss according to
3GPP TR 38.901 (a per-link stochastic LOS/NLOS draw and log-normal shadow
fading) for the terrestrial link, and free-space path loss (FSPL) under a
minimum-elevation visibility mask for the satellite link, and derives
per-user key performance indicators (KPIs): Shannon capacity and energy
per bit, the latter based on a linear power-consumption model (the EARTH
model for the base station, a power-amplifier efficiency model for the
satellite).

Beyond the static reference scenario, the simulation is extended in two
directions: repeated runs of the same topology, redrawing the channel
independently each time, to characterize KPI variance; and a simulation
of multiple consecutive transmissions across a satellite pass, with the
sub-satellite point moving at its real ground-track speed, derived from
the Keplerian orbital period, to capture how elevation, channel quality,
and handovers evolve over time. A first Monte Carlo framework further
generates a labeled dataset over randomized topologies, used to train two
classifiers (Logistic Regression, Random Forest) in three variants with
progressively less, more realistic SNR information available: exact,
instantaneous SNR (unrealistic, useful only as a pipeline sanity check),
noisy SNR standing in for a real, imperfect measurement, and no SNR at
all (geometry only).

Results show a clear split between terrestrial and satellite service
across the three KPIs examined: the satellite offers lower throughput but
markedly better energy efficiency per bit. The pass simulation initially
exposed frequent, unrealistic node switching caused by independent
(i.i.d.) per-step shadow fading; this was subsequently addressed by
adopting a spatially correlated shadow-fading model (Gudmundson, 1991),
which brought handovers down to a physically plausible level (zero
genuine handovers in the final pass scenario). The initial lack of an
explicit outage state was also addressed: when no candidate offers a
practically usable SNR, the user is now explicitly recorded as out of
coverage instead of being assigned to whichever one is "least bad." The
same usable-SNR threshold was then applied independently to each
candidate link rather than to a single winner: the user now connects
simultaneously to a terrestrial base station \emph{and} a satellite
whenever both clear the threshold (multi-connectivity, the four-state
SS-SBS model), with additive capacity/energy across both active links.
This change surfaced an unexpected finding: given this scenario's
coverage geometry, the "terrestrial-only" category nearly disappeared
from the dataset, which in turn noticeably reshaped all three machine
learning experiments below. Hysteresis and a time-to-trigger were also
added to the connectivity decision, following the 3GPP TS~38.331 Event~A3
pattern, and validated on a purpose-built stress-test scenario ($67\%$
fewer artificial state transitions). A load-based
fairness mechanism was then added, in two variants: a per-base-station
one, where an overloaded base station preferentially disconnects users
who also have a satellite alternative, protecting users with no
alternative at all -- with a measured, honestly reported benefit
($+33\%$ capacity for the protected group) and a corresponding cost for
the other; and a network-wide one, which instead looks for the most
overloaded base station across the whole network and prefers moving a
user sideways to another, less-loaded base station over dropping it to
the satellite alone -- letting even users with no satellite alternative
be actively helped, not just protected. On a purpose-built two-overlapping-cell
scenario the network-wide variant fully rebalances an overload the
per-station one cannot touch at all, though again honestly, its
greedy destination choice can dilute a previously lightly-loaded
neighbour enough to lower the population's mean capacity even as every
individually moved user improves. A full joint multi-node optimization
-- one that weighs aggregate throughput against load balance, rather
than simply chasing free capacity -- now sets the immediate next
priority for the work. For machine learning, the three
experiments spanning exact, noisy, and no SNR at all initially collapsed
to near-identical accuracy after multi-connectivity was added
($99.72\%\rightarrow99.15\%\rightarrow99.15\%$, Random Forest) --
traced to the scenario generator of the time, which never placed a user
near the actual usable-SNR boundary and had stopped producing any
terrestrial-only case at all, leaving nothing for a noisy or missing SNR
reading to actually help resolve. Fixing the generator (genuine
terrestrial-only cases, near-user placement spanning the real coverage
edge) restored a real, honestly measured gap: exact SNR reaches
$100.00\%$ accuracy, noisy SNR $93.12\%$ -- the result this thesis
treats as representative of a real deployment -- and geometry alone
$93.93\%$, a clear step below exact SNR rather than a near-tie.

\vspace{1em}
\noindent\textbf{Keywords:} multi-connectivity, satellite-terrestrial
integrated networks, 6G, LEO satellites, 3GPP TR 38.901, machine
learning, energy efficiency, handover

\end{english}
